\section{Discussion}
\label{sec:discussion}

\subsection{Vertical placement of the canopy drag}

Of the two processes examined, the seagrass canopy controls transport in this
basin, with an effect about five times that of wave coupling and the same under
both wave treatments. The mechanism the model offers is vertical rather than
frictional, since introducing the canopy removes momentum from the lower part
of the water column without removing a proportionate amount from the upper
part, more than doubling the ratio of surface to bed speed. Such a profile is
the signature of a canopy shear layer rather than of a rougher bed, which acts
through a logarithmic profile whose surface-to-bed ratio is bounded well below
the value obtained here \citep{Nepf2012}, and the bed shear stress confirms the
distinction by its sign, falling under the canopy where an equivalent bed
roughness raises it over the same cells. A surface drifter samples the top of
that profile, so a treatment that reshapes the profile relocates the drifter
even where it changes the depth-averaged transport little.

The result stands in a complementary relation to the two earlier models of this
basin. In the depth-averaged framework of \citet{Ingrassia2024} an equivalent
bed roughness is the appropriate representation of a canopy, since there is no
vertical profile for the canopy to shape, and the three-dimensional circulation
established by \citet{DeMarchis2012} is what motivates a three-dimensional
solver here. The present work quantifies what a canopy formulation resolving
the drag over the vertical extent of the plants adds in a basin of this depth.
Its parameters were taken from published measurement and not adjusted against
any observation, so the improvement follows from where the drag is applied
rather than from how it was tuned.

A start from a uniform offshore column lowers the skill of the bare members by
0.08 and 0.09 and that of the vegetated members by at most 0.02
(Section~\ref{sec:res_sensitivity}), so the canopy contrast grows from about
0.12 to about 0.19 when the start is wrong. An error elsewhere in the model
thus degrades the bare configurations and leaves the vegetated ones nearly
unaffected, and the magnitude of a canopy contrast depends on what else the
model gets wrong, even where its sign does not.

\subsection{Information carried by each gauge position}

Water level separated the members by the canopy at AltaVilaEst, by wave
coupling at the North Mouth and by neither at the South Mouth. In a micro-tidal
basin with two open connections, agreement at the inlet gauges establishes the
offshore forcing, the datum and the exchange through the mouths, the wave
contribution at the North Mouth being the one dynamic effect that reaches them,
whereas the level in the interior is set by what the water does on its way in,
which is where the canopy acts. Water level is therefore not inherently
insensitive to the canopy, since the one gauge sited inside the vegetated basin
identified the same factor as the drifters and ranked the members the same
way. Together with the drifter result, the gauge result refines the argument that water
level constrains transport weakly and that Lagrangian observations form a
distinct constraint \citep{Liu2011, Revelard2021}. For a basin of this type the
question is less whether to validate on water level than where to site the
instrument meant to discriminate, and the interior of the vegetated compartment
is a more informative position for that purpose than the inlet.

\subsection{Limitations}
\label{sec:disc_limits}

\textbf{Wave dissipation by the canopy.} In the four members of the factorial
the canopy does not reach the wave model, because the coupling passes a bed
roughness field to the spectral model while the canopy enters the momentum
equation of the flow. Seagrass meadows attenuate waves, in this species and at
this scale \citep{Infantes2012, Prinos2010}, and \texttt{WV}$^{\mathrm{s}}$
shows that the omission accounts for most of the wave penalty under the canopy
and that including it enlarges the canopy contrast with waves, from $+0.114$ to
$+0.137$. That member is a single configuration, built through an adjustment of
the wave input described in Section~S2, and the profile diagnostics of
Section~\ref{sec:res_3d} were not repeated for it.

\textbf{Canopy reconfiguration.} Three one-day runs restarted from the
stratified state of 8 July and differing only in leaf density set the bending
response of the canopy. With leaves of 900~kg~m$^{-3}$ the depth-averaged speed
over the meadow is 8 to 13\% higher than under a rigid canopy of the same
height, and leaves of 1020~kg~m$^{-3}$ bend further and leave the
900~kg~m$^{-3}$ run 13 to 20\% slower, since water over the meadow close to
1028~kg~m$^{-3}$ leaves the denser leaf about a sixteenth of the buoyancy of
the lighter one. A 13\% change in a value taken from a published table thus
moves the velocity field by as much or more. The flume measurements of
\citet{CiraoloFerreri2006} describe reconfiguration through a friction factor,
available in this solver only in its depth-averaged mode, and the two were not
compared.

\textbf{Fixed bed.} Waves govern the sediment dynamics of comparably shallow
micro-tidal basins \citep{Carniello2011}, but no repeated bathymetric survey or
in-situ turbidity record is available here to constrain the sediment
parameters, so the bed is fixed in every member and morphodynamic coupling lies
outside the present design.

\textbf{Window, restart segments and initial condition.} The record covers two
days of water level and thirty-four hours of drifter observations within one
summer season, and the vegetated members are three-day restart segments, so
seasonal behaviour and storm conditions lie outside what the design addresses.
The stratified start is a single profile applied over the domain and leaves out
the horizontal structure of the CMEMS field, so the canopy contrasts may still
carry part of a start-up adjustment, whose size the uniform start bounds. The
solution also depends on the host and decomposition at the Courant limit of the
ensemble, by 0.034 in skill, which all members share.

\textbf{\emph{Cymodocea nodosa} parameters.} The \emph{C.\ nodosa} parameters
come from a flume study of plants collected elsewhere, no measurement of that
species in this lagoon being available. The species carries about 1.5\% of the
canopy frontal area in the model, which bounds the consequence.
